% =====================================================================
\section{Wi-Fi и веб-сервер}
% =====================================================================
\lessonintro
  {подключить ESP32-S3 к Wi-Fi и сделать веб-страницу для управления лампой}
  {урок 7 (функции лампы под мьютексом), урок 5 (почему ручка на ADC1)}
  {лампой можно управлять с телефона: кнопка вкл/выкл и ползунок яркости}
  {Wi-Fi-сеть на \SI{2.4}{GHz}}

\subsection{Режимы Wi-Fi}

\begin{figure}[H]
\centering
\begin{tikzpicture}[node distance=10mm]
  \begin{scope}
    \node[blk,fill=blkrf] (r) {Роутер};
    \node[blk,fill=blkcpu,below left=10mm and 0mm of r] (e) {ESP32-S3\\ \scriptsize STA};
    \node[blk,fill=blkmem,below right=10mm and 0mm of r] (p) {Телефон/ПК};
    \draw[darr,dashed] (e) -- (r); \draw[darr,dashed] (p) -- (r);
    \node[font=\sffamily\bfseries\small,above=2mm of r] {Станция (STA)};
  \end{scope}
  \begin{scope}[xshift=6cm]
    \node[blk,fill=blkcpu] (e2) {ESP32-S3\\ \scriptsize AP 192.168.4.1};
    \node[blk,fill=blkmem,below left=10mm and -2mm of e2] (p2) {Телефон};
    \node[blk,fill=blkmem,below right=10mm and -2mm of e2] (p3) {Ноутбук};
    \draw[darr,dashed] (p2) -- (e2); \draw[darr,dashed] (p3) -- (e2);
    \node[font=\sffamily\bfseries\small,above=2mm of e2] {Точка доступа (AP)};
  \end{scope}
  \begin{scope}[xshift=12cm]
    \node[blk,fill=blkrf] (r3) {Роутер};
    \node[blk,fill=blkcpu,below=10mm of r3] (e3) {ESP32-S3\\ \scriptsize AP + STA};
    \node[blk,fill=blkmem,below=6mm of e3] (p4) {Телефон};
    \draw[darr,dashed] (e3) -- (r3); \draw[darr,dashed] (p4) -- (e3);
    \node[font=\sffamily\bfseries\small,above=2mm of r3] {Смешанный};
  \end{scope}
\end{tikzpicture}
\caption{Три режима работы Wi-Fi. Лампа будет работать как станция}
\end{figure}

Стек Wi-Fi --- это несколько задач FreeRTOS, которые по умолчанию работают на ядре~0. Наши задачи
закреплены за ядром~1, поэтому сеть не мешает кнопке и экрану.

\subsection{Подключение к сети}

\codecap{Подключение к Wi-Fi}
\codeinput{l8_connect}

\begin{caution}
ESP32-S3 работает только в диапазоне \SI{2.4}{GHz}. Сети \SI{5}{GHz} он не увидит.
\end{caution}

\subsection{Как работает HTTP-запрос}

\begin{figure}[H]
\centering
\begin{tikzpicture}[y=-0.85cm]
  \node[blk,fill=blkmem,minimum width=3cm] (b) at (0,0) {Браузер};
  \node[blk,fill=blkcpu,minimum width=3cm] (s) at (9,0) {ESP32-S3\\ \scriptsize WebServer :80};
  \draw[thick,gray] (b.south) -- (0,8.2);
  \draw[thick,gray] (s.south) -- (9,8.2);
  \draw[arr,accent] (0,1.5) -- node[above,font=\ttfamily\scriptsize]{GET / HTTP/1.1} (9,1.5);
  \draw[arr,okgreen] (9,2.6) -- node[above,font=\ttfamily\scriptsize]{200 OK + HTML-страница} (0,2.6);
  \node[font=\sffamily\scriptsize,anchor=east] at (-0.2,3.7) {пользователь жмёт «Вкл/выкл»};
  \draw[arr,accent] (0,4.4) -- node[above,font=\ttfamily\scriptsize]{GET /toggle} (9,4.4);
  \fill[esp!30] (8.85,4.6) rectangle (9.15,5.6);
  \node[font=\ttfamily\scriptsize,anchor=west] at (9.3,5.1) {lampToggle()};
  \draw[arr,okgreen] (9,6.0) -- node[above,font=\ttfamily\scriptsize]{303 → /} (0,6.0);
  \draw[arr,accent] (0,7.0) -- node[above,font=\ttfamily\scriptsize]{GET /} (9,7.0);
  \draw[arr,okgreen] (9,7.9) -- node[above,font=\ttfamily\scriptsize]{200 OK (Лампа: ВКЛ)} (0,7.9);
\end{tikzpicture}
\caption{Последовательность запросов при нажатии кнопки на веб-странице}
\end{figure}

Каждой ссылке на странице соответствует обработчик на плате. После действия сервер отвечает
перенаправлением (\code{303}) на главную страницу --- браузер сразу загружает её с новым состоянием.

\subsection{Шаг проекта 8: управление через браузер}

\progress{8}

Страница строится из текущего состояния (\code{getLamp()}), а обработчики вызывают те же функции
\code{lampToggle()} и \code{lampSetLevel()}, что кнопка и ручка. Благодаря мьютексу из урока~7
не важно, из какой задачи пришла команда. Подключение к сети не ждём: лампа работает и без Wi-Fi.

\codecap{Умная лампа, шаг 8 (изменения относительно шага 7)}
\codeinput{lamp_step8}

После загрузки IP-адрес появится в заголовке экрана. Откройте его в браузере телефона, подключённого
к той же сети, и попробуйте одновременно крутить ручку и двигать ползунок.

\subsection{Если роутера нет: собственная точка доступа}

ESP32-S3 может сам создать сеть. Замените в \code{setupWeb()} две строки подключения на:

\codeinput{l8_softap}

\begin{tip}
Чтобы не запоминать IP-адрес, подключите mDNS: \code{\#include <ESPmDNS.h>} и
\code{MDNS.begin("lamp");} --- страница откроется по адресу \code{http://lamp.local/}.
\end{tip}

\begin{summary}
\begin{itemize}[leftmargin=*]
  \item Режимы Wi-Fi: станция, точка доступа, смешанный; только \SI{2.4}{GHz}.
  \item \code{WebServer}: адрес → обработчик; \code{handleClient()} вызываем постоянно.
  \item Сеть --- ещё один источник команд; лампа принимает их через те же функции, что и кнопка.
\end{itemize}
\end{summary}

\begin{tasks}
\begin{enumerate}
  \item Сделайте адрес \code{/api/state}, возвращающий JSON: \code{\{"on":true,"level":75\}}.
  \item Если за 15 секунд не удалось подключиться к роутеру, переходите в режим точки доступа.
  \item Добавьте на страницу автообновление раз в 2 секунды, чтобы видеть изменения от ручки.
\end{enumerate}
\end{tasks}
